% Propietario reproducido por unidad demostrativa; véase PROPIETARIOS_CONTINUO_UNIVERSO.json.
% Origen: XI ES CUMPLIMIENTO_EDITORIAL_20260928, sections/continuo_cinco_construcciones_tipado.tex:1--327.
% Cuerpos y pruebas conservados; sólo namespace, rutas y remisiones contextuales declaradas.
\subsection{Joint emergence of the five constructions}
\label{hmtfg:base:sec:cinco-construcciones-continuo-ch}

For \(0\leq k\leq N\), let
\[
 \rho_{N,k}:=\rho_k\circ\rho_{k+1}\circ\cdots\circ\rho_{N-1},
 \qquad \rho_{k,k}:=\operatorname{id}_{H_k},
\]
and, for a surviving history \(\zeta\in H_k\), define
\[
 \operatorname{Desc}_{k,N}(\zeta)
 :=\{x\in H_N:\rho_{N,k}(x)=\zeta\}.
\]
The finite tree \(\mathscr T_{k,N}(\zeta)\) has vertices
\[
 \coprod_{j=k}^{N}\{y\in H_j:\rho_{j,k}(y)=\zeta\}
\]
and has as edges only the TPK prolongations
\(a:y\to z\) with \(z\in\rho_j^{-1}(y)\). Thus neither its vertices nor
its prolongations come from an auxiliary tree.

For \(x\in\operatorname{Desc}_{k,N}(\zeta)\), the marked branch
\(\lambda_{\zeta,x}\) is the sequence
\[
 \zeta=x_k\longrightarrow x_{k+1}\longrightarrow\cdots
 \longrightarrow x_N=x
\]
determined by the truncations of \(x\). The common base register is
\[
\begin{aligned}
 \mathscr R_{k,N}(\zeta;x):=
 \bigl(&\mathscr T_{k,N}(\zeta),\lambda_{\zeta,x};\\
 &\bigl(\varepsilon_a,s_a,o_a,
   (\delta_{\diamond,a},r_{\diamond,a},q_{\diamond,a})
       _{\diamond\in\{\Sigma,\Pi\}},
   \operatorname{Carr}_a,C_a,t_a,\operatorname{Fr}_a,
   \operatorname{Ruta}_a,\operatorname{Mem}_a,
   \operatorname{Surv}_a\bigr)_a;\\
 &\operatorname{src},\operatorname{tgt},\circ,
   (\rho_j)_{k\leq j<N}\bigr),
\end{aligned}
\tag{C1}
\label{hmtfg:base:eq:base-comun-registros}
\]
where
\[
 U_a(m_0)=C_am_0+t_a,
 \qquad
 \delta_{\diamond,a}=r_{\diamond,a}+9q_{\diamond,a}.
\]
The symbol \(m_0\) in this formula is a memory coordinate; it is not the
modular depth \(m\) used below. The true base space is
\[
 \mathfrak R_{k,N}(\zeta)
 :=\{\mathscr R_{k,N}(\zeta;x):
       x\in\operatorname{Desc}_{k,N}(\zeta)\}.
\]
The marked branch is not an added input: it is the sequence of actually
produced edges ending at \(x\). The rest of the register simultaneously
preserves the lateral fibres of the tree; thus the mark does not replace
branching with an isolated path.

\Needspace{15\baselineskip}
Set \(R_m:=\mathbb Z/3^m\mathbb Z\), and let
\(V_j(\zeta):=\operatorname{Desc}_{k,j}(\zeta)\). Transport, the
sheet projectors and the terminal residue restrict to the subtree through
\[
 \mathcal U_j^\zeta:=
 \mathcal U_j\!\upharpoonright_{\ell^2(V_j(\zeta),\mu_j)},
 \qquad P_j^\zeta:=P_j\!\upharpoonright_{V_j(\zeta)},
 \qquad Q_j^\zeta:=I-P_j^\zeta,
\]
\[
 A_j^\zeta:=P_{j+1}^\zeta\mathcal U_j^\zeta P_j^\zeta
             +Q_{j+1}^\zeta\mathcal U_j^\zeta Q_j^\zeta,
 \qquad
 \eta_j^\zeta:=Q_{j+1}^\zeta\mathcal U_j^\zeta P_j^\zeta
             +P_{j+1}^\zeta\mathcal U_j^\zeta Q_j^\zeta.
\]
For \(k\leq n\leq N\), let
\[
 F_{k,k}^\zeta:=I,
 \qquad
 F_{k,n}^\zeta:=A_{n-1}^\zeta\cdots A_k^\zeta,
 \qquad
 T_{k,n}^\zeta:=(F_{k,n}^\zeta)^*F_{k,n}^\zeta.
\]
The ports belonging to the same subtree are
\[
 \mathscr P_j^\Sigma(\zeta):=
 \{(y,\Sigma):y\in V_j(\zeta)\},
 \qquad
 \mathscr P_j^\Pi(\zeta):=
 \{(y,\Pi):y\in V_j(\zeta)\};
\]
we denote by \(B_j^\zeta,\kappa_j^\zeta,\delta_j^\zeta\) the restrictions
of the incidence and the rectangular complex to those ports, and by
\(L_j^\zeta\) the precomposition induced by removing prefixes.

To define the path reader without abbreviations, let
\[
 W_\ell^\zeta(y):=
 R_m[\operatorname{Path}_{\mathscr T_{k,\ell}(\zeta)}
          (y,\operatorname{Term}_\ell)]
\]
and set
\[
 \mathscr D_{q,r;k,\ell}^{(m)}[\mathscr R]
 :=\bigoplus_{y_0\to\cdots\to y_q\ \text{in}\
                  \mathscr T_{k,\ell}(\zeta)}
 W_\ell^\zeta(y_q)\otimes_{R_m}
 \bigl(\Gamma_r(y_0)\otimes_{\mathbb Z}R_m\bigr),
 \qquad
 \operatorname{Tot}_d\mathscr D
 :=\bigoplus_{q+r=d}\mathscr D_{q,r;k,\ell}^{(m)}[\mathscr R].
\]
The formulas already proved produce, on each
\(\mathscr R\in\mathfrak R_{k,N}(\zeta)\), the five objects
\[
\begin{aligned}
 \mathcal F_{\rm sp}[\mathscr R]
 &:=(\ell^2(V_j(\zeta),\mu_j),
       \mathcal U_j^\zeta,A_j^\zeta,\eta_j^\zeta,
       T_{k,j}^\zeta)_{k\leq j\leq N},\\
 \mathcal F_{\rm sol}[\mathscr R]
 &:=(b,\kappa_{\rm can},C_\alpha,\kappa_\alpha)
      _{\alpha\subset\mathscr T_{k,N}(\zeta)},\\
 \mathcal F_{\rm gau}^{(m)}[\mathscr R]
 &:=(\mathscr P_j^\Sigma(\zeta),\mathscr P_j^\Pi(\zeta),
       B_j^\zeta,\kappa_j^\zeta,\delta_j^\zeta;R_m)_{k\leq j\leq N},\\
 \mathcal F_{\rm coh}^{(m)}[\mathscr R]
 &:=(\Gamma_y\otimes R_m,\partial,\Psi_\alpha,H_\bullet)
      _{y,\alpha\subset\mathscr T_{k,N}(\zeta)},\\
 \mathcal F_{\rm det}^{(m)}[\mathscr R]
 &:=(\operatorname{Det}(\operatorname{Tot}
       \mathscr D_{\bullet,\bullet;k,\ell}^{(m)}[\mathscr R]))
       _{k\leq\ell\leq N}.
\end{aligned}
\tag{C2}
\label{hmtfg:base:eq:cinco-lectores-comunes}
\]
Here \(\mathscr D_{\bullet,\bullet;k,\ell}^{(m)}[\mathscr R]\) is the normalized path-and-memory bicomplex of the preceding sections, restricted to the subtree of horizon \(\ell\) and reduced over \(R_m\). The last reader preserves the ordered collection of determinant lines up to horizon \(N\); in this way, its truncation is a canonical projection and does not presuppose a nonexistent determinantal factorization. The law \(\mu\) is the canonical census law, and the ports are fixed through the lexicographically minimal branch of the APP order. Thus none of the five readers receives external weighting or an external port.

For
\(T\in\mathfrak T:=\{{\rm sp},{\rm sol},{\rm gau},{\rm coh},{\rm det}\}\),
define the typed total space
\[
 \widetilde{\mathcal F}_{T;k,N}^{(m)}(\zeta)
 :=\coprod_{\mathscr R\in\mathfrak R_{k,N}(\zeta)}
   \{\mathcal F_T^{(m)}[\mathscr R]\}_{\mathscr R},
 \qquad
 p_T:\widetilde{\mathcal F}_{T;k,N}^{(m)}(\zeta)
      \longrightarrow\mathfrak R_{k,N}(\zeta),
\]
where
\(\mathcal F_{\rm sp}^{(m)}:=\mathcal F_{\rm sp}\) and
\(\mathcal F_{\rm sol}^{(m)}:=\mathcal F_{\rm sol}\); that is, the
superscript \(m\) is inert for those two readers.
The projection \(p_T\) forgets only the reader structure and preserves
its source register. The correct fibre product is
\[
\begin{aligned}
 \mathcal J_{k,N}^{(m)}(\zeta):={}&
 \widetilde{\mathcal F}_{\rm sp;k,N}^{(m)}(\zeta)
 \mathop{\times}_{\mathfrak R_{k,N}(\zeta)}
 \widetilde{\mathcal F}_{\rm sol;k,N}^{(m)}(\zeta)
 \mathop{\times}_{\mathfrak R_{k,N}(\zeta)}
 \widetilde{\mathcal F}_{\rm gau;k,N}^{(m)}(\zeta)\\
 &\mathop{\times}_{\mathfrak R_{k,N}(\zeta)}
 \widetilde{\mathcal F}_{\rm coh;k,N}^{(m)}(\zeta)
 \mathop{\times}_{\mathfrak R_{k,N}(\zeta)}
 \widetilde{\mathcal F}_{\rm det;k,N}^{(m)}(\zeta).
\end{aligned}
\tag{C3}
\label{hmtfg:base:eq:constructor-correlativo}
\]
Equality of the five projections requires preserving the same tree, the
same edges, the same produced branch and the same register. For each
\(\mathscr R\), the five constructions determine the element
\[
 s_{k,N}^{(m)}(\mathscr R)
 :=(\mathcal F_T^{(m)}[\mathscr R])_{T\in\mathfrak T}
 \in\mathcal J_{k,N}^{(m)}(\zeta).
\]
Existence of this element follows from the operations of
\eqref{hmtfg:base:eq:cinco-lectores-comunes}, not from the
fibre-product notation.

Removal of the last layer induces
\[
 \operatorname{res}^{\mathfrak R}_{N+1,N}
 \bigl(\mathscr R_{k,N+1}(\zeta;x)\bigr)
 :=\mathscr R_{k,N}(\zeta;\rho_N(x)).
\]
Restriction of each operator to the retained generators defines
\[
 r^m_{N+1,N}:\mathcal J_{k,N+1}^{(m)}(\zeta)
 \longrightarrow\mathcal J_{k,N}^{(m)}(\zeta).
\]
In the determinantal reader, \(r^m_{N+1,N}\) removes only the last coordinate of the collection of lines of \eqref{hmtfg:base:eq:cinco-lectores-comunes}. The reduction \(R_{m+1}\twoheadrightarrow R_m\) defines, by base change,
\[
 q^N_{m+1,m}:\mathcal J_{k,N}^{(m+1)}(\zeta)
 \longrightarrow\mathcal J_{k,N}^{(m)}(\zeta).
\]
On each generator, we have
\[
\begin{gathered}
 r^m_{N+1,N}q^{N+1}_{m+1,m}
 =q^N_{m+1,m}r^{m+1}_{N+1,N},\\
 \partial\Psi_\alpha=\Psi_\alpha\partial,
 \qquad L_j^\zeta\kappa_j^\zeta
       =\kappa_{j+1}^\zeta L_j^\zeta,
 \qquad L_j^\zeta\delta_j^\zeta
       =\delta_{j+1}^\zeta L_j^\zeta,
\end{gathered}
\tag{C4}
\label{hmtfg:base:eq:naturalidad-D}
\]
and compositions of the maps \(r\) and \(q\) satisfy the functorial
identities. The first equality holds because deleting the last layer and
reducing the coefficients act on distinct coordinates. The other three
are the naturality identities proved earlier. A unique
pro-system in \((N,m)\) is obtained, not five juxtaposed limits.

At the root, define
\[
 \mathfrak J_{N,\bullet}^{(m)}:=\mathcal J_{0,N}^{(m)}(\ast).
\tag{C5}
\label{hmtfg:base:eq:fibras-conjuntas-horizonte}
\]
The joint discrete structure of the continuum is the limit of those
correlative objects:
\[
 \boxed{
 \mathcal C_{\rm cont}^{\rm disc}
 :=\varprojlim_{\substack{N\geq1\\m\geq1}}
   \bigl(\mathfrak J_{N,\bullet}^{(m)},
         r^m_{N+1,N},q^N_{m+1,m}\bigr).}
\tag{C6}
\label{hmtfg:base:eq:continuo-discreto}
\]
This is a definition through the pro-system generated by the complete tree; it is not
the image or graph of a map retaining the input state as
its first coordinate.

The registers have a produced terminal mark. Their compatible projections
induce
\[
 \Pi:\mathcal C_{\rm cont}^{\rm disc}\longrightarrow H_\infty^{\rm enr}.
\]
Conversely, a complete history determines, through its already emitted prefixes,
the compatible sections \(s_{0,N}^{(m)}\). Thus, after having
defined the continuum, one obtains
\[
 \operatorname{Gen}_{\rm cont}:H_\infty^{\rm enr}
 \longrightarrow\mathcal C_{\rm cont}^{\rm disc},
 \qquad
 \Pi\circ\operatorname{Gen}_{\rm cont}
 =\operatorname{id}_{H_\infty^{\rm enr}}.
\tag{C7}
\label{hmtfg:base:eq:seccion-historias-continuo}
\]
This section proves recoverability; it is not used to define
\(\mathcal C_{\rm cont}^{\rm disc}\).

\begin{theorem}[Correlative emergence of the five constructions]
\label{hmtfg:base:thm:correlacion-cinco}
For every surviving history \(\zeta\in H_k\), every \(N>k\), every
\(m\geq1\) and every \(\mathscr R\in\mathfrak R_{k,N}(\zeta)\), the five
constructions of \eqref{hmtfg:base:eq:cinco-lectores-comunes} are defined on the
same register, and \eqref{hmtfg:base:eq:constructor-correlativo} contains
\(s_{k,N}^{(m)}(\mathscr R)\). Their truncations, modular reductions and
transports satisfy \eqref{hmtfg:base:eq:naturalidad-D}. Furthermore,
\(\mathcal C_{\rm cont}^{\rm disc}\neq\varnothing\), the map \(\Pi\)
is surjective and is split by \(\operatorname{Gen}_{\rm cont}\). The
five components are inseparable aspects of the same continuum, not five
domains or five independent maps.
\end{theorem}

\begin{proof}
The survival defining \(H_k\) and the TPK prolongation law give
\(\rho_k^{-1}(\zeta)\neq\varnothing\) for each vertex. Since the fibres are
finite, Kőnig’s lemma produces \(H_\infty^{\rm enr}\neq\varnothing\).

Given a register \(\mathscr R\), the transport of histories and sheets produces
\(\mathcal F_{\rm sp}[\mathscr R]\); affine composition and the exact cancellation
defect produce \(\mathcal F_{\rm sol}[\mathscr R]\); incidences
and the rectangular sequence produce \(\mathcal F_{\rm gau}^{(m)}[\mathscr R]\);
the memory complex and its transport produce
\(\mathcal F_{\rm coh}^{(m)}[\mathscr R]\); and the lines of the normalized
path assembly produce \(\mathcal F_{\rm det}^{(m)}[\mathscr R]\).
Each construction acts on the same generated edges, not on an
assumed subdomain or on the abstract ambient space of \(729\) words.

The affine, transport, incidence and memory compatibilities prove \eqref{hmtfg:base:eq:naturalidad-D}. Truncating the determinantal collection of \eqref{hmtfg:base:eq:cinco-lectores-comunes} commutes by definition with \(r\); its base change commutes with \(q\). The sections \(s_{0,N}^{(m)}\) are compatible, so they produce \(\operatorname{Gen}_{\rm cont}\). The identity \eqref{hmtfg:base:eq:seccion-historias-continuo} simultaneously proves nonemptiness of the limit, surjectivity of \(\Pi\) and recoverability of each history without introducing it as a defining coordinate.

The fibre, relation, number, orientation, carry, memory,
boundary and route remain in the common register; the multisection,
terminals, and linear and probabilistic readers remain in the five
correlative constructions. Passing to the limit therefore preserves dependency
closure and does not separate the continuum.
\end{proof}

Subsequent recognition of the uniquely constructed continuum through spectral,
solenoidal, gauge, cohomological or determinantal language belongs to the
later conventional comparison. It does not select its states or enter
its APP--TRIT--TPK generation.
